%---------------------------------------------------------------------
\chapter{Phase 2 --- The Fourteen-Lane Evaluation Campaign}
\label{ch:campaign}

Phase 2's second axis was empirical scale. Phase 1 evaluated a diagnostic harness
on a small number of studies. Phase 2 evaluated a single trained actor against
\textbf{fourteen held-out agent and RL benchmark suites} under strict governance.
This chapter specifies the evaluated model, the governance that constrains what
may be claimed, and the scope semantics that keep the results comparable. The
results themselves are in Chapters~\ref{ch:p2results} and the appendices.

\section{The evaluated actor}

The campaign evaluates \textbf{one} actor under a single frozen configuration:

\begin{table}[h]
  \centering
  \caption{Evaluated actor and adapter, as recorded in the campaign receipts.}
  \label{tab:actor}
  \small
  \begin{tabular}{@{}L{0.26\linewidth}L{0.68\linewidth}@{}}
    \toprule
    \textbf{Component} & \textbf{Value} \\
    \midrule
    Base model & \texttt{Qwen/Qwen3.6-35B-A3B}, commit \texttt{995ad96e\ldots} \\
    Adapter & \texttt{arvindcr4/pavlov-portfolio-qwen36-seed809-stepfinal-tinker-cf0ad8c1\ldots}, commit \texttt{64444133\ldots} \\
    Precision & bfloat16 (merged weights for LAB-Bench, AgentDojo, Omni-MATH) \\
    Training back-end & Tinker (managed GRPO) \\
    Sampling & original Tinker sampler for VerilogEval; merged BF16 serving for the rest \\
    Math grading & separate Omni-Judge grader, later grading campaign on Modal A100 \\
    \bottomrule
  \end{tabular}
\end{table}

Because the same actor, task, reward function, and decoding configuration are held
fixed across suites, observed differences are attributable to the \emph{stack}
rather than to the experiment --- the attribution commitment inherited from
Phase~1. Two serving back-ends were used (eager-serving, graph-serving at 32{,}768
and 65{,}536-token context). \textbf{No numerical parity between serving back-ends
is claimed.} A fresh metadata request for the original sampler returns checkpoint
not found, and that gap is recorded rather than papered over.

\section{Governance: receipts, budget, decontamination}

Every lane is governed by four controls, applied uniformly:

\begin{enumerate}[leftmargin=1.6em,itemsep=3pt]
  \item \textbf{Receipt discipline.} A lane may advance only on a surviving
  receipt: a sealed request, a validation record, or a native evaluation log.
  Sealed artefacts carry SHA-256 hashes; the verifier re-derives them.
  \item \textbf{Explicit budget attribution.} Managed compute is charged per lane.
  Spend is tracked against an authorised cap and reported as a ledger, never
  blended across lanes.
  \item \textbf{Decontamination.} Source-overlap checks are run against a defined
  source envelope. This is a bounded check: it covers the declared envelope, not
  images, semantic similarity, or universal pretraining exposure. The bound is
  stated wherever the check is reported.
  \item \textbf{Scope separation.} Where a suite has an \emph{original contract}
  and a later \emph{replacement scope}, the two are never pooled. A replacement
  scope is named as such in every table.
\end{enumerate}

\section{Scope semantics: original contract vs replacement scope}

The single most important reporting decision in this campaign is the separation of
scope. Each of the fourteen lanes is in exactly one of five states, and the state
--- not a percentage alone --- is what determines whether a number may be quoted:

\begin{table}[h]
  \centering
  \caption{Lane state vocabulary. A percentage without its state is not a
  result.}
  \label{tab:states}
  \small
  \begin{tabular}{@{}L{0.20\linewidth}L{0.74\linewidth}@{}}
    \toprule
    \textbf{State} & \textbf{Meaning} \\
    \midrule
    COMPLETE & Every item in the declared scope carries a native, verified verdict. A score may be quoted with its denominator. \\
    TERMINAL & All dispositions are recorded, including explicit skip reasons, under the native protocol. Coverage is quoted; score semantics are stated separately. \\
    PARTIAL & A verified score exists over a stated sub-scope. The score is quoted \emph{with its coverage}; the full-suite score is explicitly unavailable. \\
    BLOCKED & No local or paid path to the artefacts exists (private bundle, missing payload, unanswered access request). No score is quoted. \\
    NOT RUN & Preconditions are unmet within budget or quota. No score is quoted. \\
    \bottomrule
  \end{tabular}
\end{table}

Three replacement scopes closed strictly \textbf{COMPLETE} (E8 public, E10
benign, E11 native). One closed \textbf{TERMINAL} at 99.95\% coverage (E14
Omni-MATH). The original-contract \textbf{E1} and \textbf{E11} suites carry full
verified scores. The remainder are recorded as PARTIAL or BLOCKED with immutable
evidence trails. Seven lanes have \emph{no} local or paid path at all, because
the required artefacts are private, held out, or gated behind an unanswered
access request:

\begin{itemize}[leftmargin=1.4em,itemsep=2pt]
  \item E3 --- SDAB private bundle
  \item E7 --- BinaryAudit private payload
  \item E8-original --- LifeSciBench package
  \item E10-original --- AgentHarm private tasks
  \item E12 --- AppBench deployment
  \item E13-original --- OpenReward held-out games
  \item E14-original --- FrontierMath hosted evaluation
\end{itemize}

Access requests were sent for all seven. \textbf{None answered.} This is reported
as a finding about artefact accessibility in this benchmark ecosystem, not as a
gap in effort.

\section{Why a blocked lane is a result}

It is worth stating explicitly, because it is the question every reviewer of a
multi-lane campaign asks first: why are seven lanes not done? The answer is a
property of the ecosystem, and it is measurable.

\begin{itemize}[leftmargin=1.4em,itemsep=2pt]
  \item The suites in question gate their task payloads behind private bundles,
  hosted evaluation services, or per-institution deployment. There is no public
  download and no paid acquisition path.
  \item The campaign documented the dependency for each lane \emph{before}
  attempting it, so the blocked state is a predicted outcome rather than a
  discovered failure.
  \item Requests were filed, tracked, and left unanswered. Each unanswered request
  is itself a receipt.
\end{itemize}

A rigorous multi-suite campaign therefore reports \emph{coverage of the
accessible set} and \emph{accessibility of the inaccessible set}. Conflating the
two --- by quoting a partial score as if it were the suite score, or by silently
omitting the blocked lanes --- is the failure mode this campaign was designed to
avoid. Chapter~\ref{ch:rigour} returns to this as a threat to validity.

\section{Evidence accounting and loss}

Full transparency requires recording what happened to the evidence, not only to
the runs. One execution-source loss occurred and is retained in the record:

\begin{itemize}[leftmargin=1.4em,itemsep=2pt]
  \item A \texttt{.codex-run} deletion removed the finish-era recovery and build
  code --- specifically the E1 wave10 controller, E9 restored dependencies, and
  the E6 validator.
  \item The loss is \textbf{confined to recovery/build code}. The replacement-lane
  native runners survive intact under the flagship program directories.
  \item Sealed requests, validation receipts, and test logs pin any faithful
  re-implementation, so the loss degrades \emph{convenience} of reproduction, not
  the validity of any reported score.
  \item The verifier carries an explicit check for this condition
  (\texttt{codex\_run\_execution\_sources\_absent}), which asserts the sources are
  absent rather than pretending they are present.
\end{itemize}

The 2026-09-12 halt was clean: no orphan processes, deployments, or sessions
remained, and disk was recovered. Freezing the campaign at a known-clean boundary
is what makes the terminal states in Chapter~\ref{ch:p2results} trustworthy.

\section{Terminal-state directive and decision package}

The campaign was closed by an explicit directive rather than by attrition. A
decision package enumerated six user-level decisions (E2/E13 lifecycle amendment;
E5 successor budget; E1 re-implementation and incremental budget; E4 cap and
scope; AWS quota requests; external follow-ups), each with a recommendation
derived from the research-value ordering described in Chapter~\ref{ch:rigour}.
The terminal directive was recorded verbatim at the time and treated as
acceptance of those recommendations. The per-lane ledger of record is the
terminal-state table, and the actions taken under the directive are:

\begin{itemize}[leftmargin=1.4em,itemsep=2pt]
  \item \textbf{E14 closed terminal-complete.} The two parser-exclusion rows were
  resolved by reading the native judge outputs directly: both are truncated before
  any ``Equivalence Judgement'' section, so the judge emitted no verdict, and the
  native scorer's omission of missing-verdict rows is \emph{correct official
  behaviour}, not a parser bug. Score impact is none --- official accuracy
  $2271/4428 = 0.5131043831902395$ already counts both rows in the denominator as
  not-correct. The $4426/4428$ figure is \emph{parser coverage}, not the score
  denominator. No re-judge was performed: it would deviate from the official
  scorer and generate new evidence with no score effect.
  \item Remaining lanes held at their documented terminal states, with budget and
  quota blockers recorded rather than worked around.
\end{itemize}

This distinction --- coverage versus score denominator --- is the kind of
precision the campaign insists on, and it is why E14 is reported as 99.95\%
coverage at 51.31\% official accuracy rather than as a single misleading
percentage.
